% GENERATED FILE. Edit canonical lessons, then run npm run generate:books.
% canonical-source-hash: fnv1a64:82b522f93cc80f55
% canonical-lessons: ZH-C103-yong, ZH-W103-yong, ZH-C103-na, ZH-C103-reng, ZH-C103-xiu, ZH-C103-huai

\chapter{Using Things}
\label{ch:zh-using-things}
\hlchaptermodality{103}

\begin{chapteropening}
\textbf{By the end of this chapter:} \emph{I can say use, take, throw away and repair, and say something is broken.}

Complete the last lesson of chapter 103: I can say use, take, throw away and repair, and say something is broken.
\end{chapteropening}

\section[\texorpdfstring{\zh{用} (yòng)}{用 (yòng)}]{\zh{用} (yòng) — to use}

\label{lesson:ZH-C103-yong}



\begin{quote}
Before the new one: say the Mandarin for beer, then the Mandarin for biscuit.
\end{quote}

\subsection*{You'll want to know: \zh{用}}
\textbf{\zh{用}} — \emph{yòng} — \textquotedblleft{}to use\textquotedblright{}.

\begin{quote}
\textbf{\zh{我用水洗菜}} \emph{wǒ yòng shuǐ xǐ cài} — I wash the vegetables in water.
\end{quote}

The new character, in pieces: \zh{用} — the frame of the moon piece in \zh{明}, open at the foot, with two bars inside and a vertical down the middle. The next lesson puts it in your hand.

\begin{cousinweb}
\textbf{\zh{用}} (\emph{yòng}) is inherited from Middle Chinese \emph{yowngH} (Guangyun, in Baxter's transcription), glossed 'to employ' and 'by means of'.
\end{cousinweb}

\begin{grammarlens}[title={What you've built}]
One more word for everyday actions.
\end{grammarlens}

\subsection*{Your turn}
\begin{itemize}
\raggedright
  \item \emph{Say it:} \emph{yòng}
  \item \emph{Say it:} \emph{yòng}, once more
  \item \emph{Say it:} \emph{wǒ yòng shuǐ xǐ cài} — I wash the vegetables in water.
  \item \emph{Recall:} say the Mandarin for beer, then the Mandarin for biscuit, then say \emph{yòng} again
\end{itemize}

\begin{tcolorbox}[breakable,colback=teal!4,colframe=teal!35!black,title={Before you move on}]
What does \zh{用} mean? (\textquotedblleft{}To use\textquotedblright{}.) Say it once more.
\end{tcolorbox}
\section[\texorpdfstring{\zh{用} (yòng)}{用 (yòng)}]{\zh{用} — the moon frame, two bars, a vertical}

\label{lesson:ZH-W103-yong}



\begin{quote}
The word first, then the character inside it.

\begin{itemize}
\raggedright
  \item \emph{Recall:} say \emph{yòng} — \textbf{R1}, one lesson back, before it decays
\end{itemize}

\begin{quote}
\zh{明} \zh{能}
\end{quote}

\emph{Read it:} each one aloud, with nothing said first — \emph{míng}, \emph{néng}
\end{quote}

\subsection*{Script — the moon frame, two bars, a vertical}
\hlblockfigure{\detokenize{figures/ZH-W103-yong-filmstrip.pdf}}{How \zh{用} is written, stroke by stroke}

\begin{quote}
\zh{用}
\end{quote}

\textbf{The frame, two strokes}, as on the right of \zh{明}: a left-falling stroke down the left side, then a stroke across the top that turns down the right side and hooks.

\textbf{Two horizontals inside}, one above the other.

\textbf{A vertical down the middle}, from the top through both bars to the foot.

2 + 2 + 1 = five.

\subsection*{Your turn}
\begin{itemize}
\raggedright
  \item \emph{Write it:} \zh{用}, all five, with the model in front of you
  \item \emph{Cover:} write it again from memory, and count the pieces, not the strokes
\end{itemize}

\begin{tcolorbox}[breakable,colback=teal!4,colframe=teal!35!black,title={Before you move on}]
How many strokes in \zh{用}? (\textbf{Five}.) Which character you know has the same frame? (\textbf{\zh{明}}, on its right.) Which stroke comes last? (\textbf{The vertical down the middle}.)

Source: \href{https://raw.githubusercontent.com/chanind/hanzi-writer-data/68d10a4b21150cae5e1ebbd223eed289cf32d90c/data/\%E7\%94\%A8.json}{Hanzi Writer Data \zh{用}, pinned stroke paths and medians}.
\end{tcolorbox}
\section[ná]{ná — to take; to hold}

\label{lesson:ZH-C103-na}



\begin{quote}
Before the new one: say the Mandarin for biscuit, then the Mandarin for to use.
\end{quote}

\subsection*{You'll want to know: ná}
\emph{ná} — \textquotedblleft{}to take; to hold\textquotedblright{}.

\begin{quote}
\emph{nǐ ná wǒ de yàoshi le ma} — Did you take my keys?
\end{quote}

Its characters come later; for now it is a word for the ear.

\begin{cousinweb}
\emph{ná} is probably inherited from Middle Chinese \emph{nrae} (Guangyun, in Baxter's transcription), which the dictionary lists under an older form of the character, glossed 'to pull'. That the two are one word is likely; how 'to pull' became 'to take' is not settled.
\end{cousinweb}

\begin{grammarlens}[title={What you've built}]
One more word for everyday actions.
\end{grammarlens}

\subsection*{Your turn}
\begin{itemize}
\raggedright
  \item \emph{Say it:} \emph{ná}
  \item \emph{Say it:} \emph{ná}, once more
  \item \emph{Say it:} \emph{nǐ ná wǒ de yàoshi le ma} — Did you take my keys?
  \item \emph{Recall:} say the Mandarin for biscuit, then the Mandarin for to use, then say \emph{ná} again
\end{itemize}

\begin{tcolorbox}[breakable,colback=teal!4,colframe=teal!35!black,title={Before you move on}]
What does ná mean? (\textquotedblleft{}To take; to hold\textquotedblright{}.) Say it once more.
\end{tcolorbox}
\section[rēng]{rēng — to throw; to throw away}

\label{lesson:ZH-C103-reng}



\begin{quote}
Before the new one: say the Mandarin for to use, then the Mandarin for to take hold of.
\end{quote}

\subsection*{You'll want to know: rēng}
\emph{rēng} — \textquotedblleft{}to throw; to throw away\textquotedblright{}.

\begin{quote}
\emph{bú yào rēng wǒ de shū} — Don't throw my books away.
\end{quote}

Its characters come later; for now it is a word for the ear.

\begin{cousinweb}
\emph{rēng} is inherited from Middle Chinese \emph{nying} (Guangyun, in Baxter's transcription), glossed 'to pull'. The regular change from \emph{nying} gives réng, with a rising tone, not the level rēng of today, and throwing is not given there; how the word came to its sound and sense is not settled.
\end{cousinweb}

\begin{grammarlens}[title={What you've built}]
One more word for everyday actions.
\end{grammarlens}

\subsection*{Your turn}
\begin{itemize}
\raggedright
  \item \emph{Say it:} \emph{rēng}
  \item \emph{Say it:} \emph{rēng}, once more
  \item \emph{Say it:} \emph{bú yào rēng wǒ de shū} — Don't throw my books away.
  \item \emph{Recall:} say the Mandarin for to use, then the Mandarin for to take hold of, then say \emph{rēng} again
\end{itemize}

\begin{tcolorbox}[breakable,colback=teal!4,colframe=teal!35!black,title={Before you move on}]
What does rēng mean? (\textquotedblleft{}To throw; to throw away\textquotedblright{}.) Say it once more.
\end{tcolorbox}
\section[xiū]{xiū — to repair; to fix}

\label{lesson:ZH-C103-xiu}



\begin{quote}
Before the new one: say the Mandarin for to take hold of, then the Mandarin for to throw away.
\end{quote}

\subsection*{You'll want to know: xiū}
\emph{xiū} — \textquotedblleft{}to repair; to fix\textquotedblright{}.

\begin{quote}
\emph{tā huì xiū diànnǎo} — He knows how to fix computers.
\end{quote}

Its characters come later; for now it is a word for the ear.

\begin{cousinweb}
\emph{xiū} is inherited from Middle Chinese \emph{sjuw} (Guangyun, in Baxter's transcription), glossed 'to put in order' and 'to adorn'.
\end{cousinweb}

\begin{grammarlens}[title={What you've built}]
One more word for everyday actions.
\end{grammarlens}

\subsection*{Your turn}
\begin{itemize}
\raggedright
  \item \emph{Say it:} \emph{xiū}
  \item \emph{Say it:} \emph{xiū}, once more
  \item \emph{Say it:} \emph{tā huì xiū diànnǎo} — He knows how to fix computers.
  \item \emph{Recall:} say the Mandarin for to take hold of, then the Mandarin for to throw away, then say \emph{xiū} again
\end{itemize}

\begin{tcolorbox}[breakable,colback=teal!4,colframe=teal!35!black,title={Before you move on}]
What does xiū mean? (\textquotedblleft{}To repair; to fix\textquotedblright{}.) Say it once more.
\end{tcolorbox}
\section[huài]{huài — broken; bad}

\label{lesson:ZH-C103-huai}



\begin{quote}
Before the new one: say the Mandarin for to throw away, then the Mandarin for to repair.
\end{quote}

\subsection*{You'll want to know: huài}
\emph{huài} — \textquotedblleft{}broken; bad\textquotedblright{}.

\begin{quote}
\emph{wǒ de shǒujī huài le} — My phone's broken.
\end{quote}

Its characters come later; for now it is a word for the ear.

\begin{cousinweb}
\emph{huài} is inherited from Middle Chinese \emph{hweajH} (Guangyun, in Baxter's transcription), glossed 'to break of itself'.
\end{cousinweb}

\begin{grammarlens}[title={What you've built}]
One more word for everyday actions.
\end{grammarlens}

\subsection*{Your turn}
\begin{itemize}
\raggedright
  \item \emph{Say it:} \emph{huài}
  \item \emph{Say it:} \emph{huài}, once more
  \item \emph{Say it:} \emph{wǒ de shǒujī huài le} — My phone's broken.
  \item \emph{Recall:} say the Mandarin for to throw away, then the Mandarin for to repair, then say \emph{huài} again
\end{itemize}

\begin{tcolorbox}[breakable,colback=teal!4,colframe=teal!35!black,title={Before you move on}]
What does huài mean? (\textquotedblleft{}Broken; bad\textquotedblright{}.) Say it once more.
\end{tcolorbox}
